% fluids_bubbles: two soap bubbles sharing a film (Vol II Ch 2, fig:fluids-bubbles)
% r1 = 2.0 cm, r2 = 3.0 cm.  Excess pressures 4S/r1 > 4S/r2, so the shared film
% carries the difference: 4S/R = 4S/r1 - 4S/r2, R = r1 r2/(r2 - r1) = 6.0 cm,
% bulging into the larger bubble.  Geometry (computed in python, units of cm):
% the three films meet at 120 deg, which puts the centres d = sqrt(r1^2 + r2^2 - r1 r2)
% = 2.646 apart; the junctions at x = 0.378, y = +-1.964 (centre 1 at the origin);
% the shared film's centre at x = -5.292.  Arc angles: bubble 1 from 79.1 to 280.9 deg,
% bubble 2 from -139.1 to 139.1 deg, shared film from -19.1 to 19.1 deg.
% Drawn at 0.55 cm per cm.
\begin{tikzpicture}[x=0.55cm,y=0.55cm,
    film/.style={draw=cPlum,construction},
    inside/.style={fill=cPlum!8,draw=none}]
  \def\d{2.6458}\def\xj{0.3780}\def\yj{1.9640}\def\xc{-5.2915}
  % interiors (pale), then the films
  \fill[inside] (0,0) circle[radius=2];
  \fill[inside] (\d,0) circle[radius=3];
  \draw[film] (\xj,\yj) arc[start angle=79.107,end angle=280.893,radius=2];
  \draw[film] (\xj,-\yj) arc[start angle=-139.107,end angle=139.107,radius=3];
  \draw[film] (\xj,-\yj) arc[start angle=-19.107,end angle=19.107,radius=6];
  % centres and radii
  \node[point mass] at (0,0) {};
  \node[point mass] at (\d,0) {};
  \draw[dim] (0,0) -- (235:2) node[midway,anchor=north west,cSteel,inner sep=1pt] {$r_1$};
  \draw[dim] (\d,0) -- ++(-35:3) node[midway,anchor=south west,cSteel,inner sep=1pt] {$r_2$};
  % labels
  \node[align=center,font=\footnotesize] at (-0.4,0.85) {higher\\ pressure};
  \node[align=center,font=\footnotesize] at (3.6,1.0) {lower\\ pressure};
  \draw[leader] (0.72,-0.9) -- (1.6,-3.6) node[anchor=north,cPlum,align=center,font=\footnotesize]
     {shared film, radius $R$,\\ bulging into the larger bubble};
\end{tikzpicture}
